\section{Introduction}
\label{sec:intro}

% Vision--Language--Action (VLA) models have recently emerged as a promising paradigm for robotic manipulation, leveraging large pre-trained Vision--Language Models (VLMs) to map rich visual observations and language instructions directly to actions~\cite{kim2025openvla,black2410pi0}. By inheriting semantic and visual priors from internet-scale data, these models exhibit strong generalization across diverse tasks and environments. However, most VLAs reason primarily in 2D image space and generate actions reactively from the current observation, without anticipating how robot motion will alter the underlying 3D scene. This 2D-centric and myopic formulation limits performance in manipulation tasks that require precise spatial reasoning, object-centric geometry understanding, and physically consistent long-horizon behavior.

% A natural direction is to augment policies with a \emph{world model} capable of predicting future environment states~\cite{lu2025gwm,cen2025worldvla}. Early world models emphasized latent dynamics for planning or reinforcement learning, but these lacked the explicit geometry needed for contact-rich manipulation. More recent approaches predict future RGB frames or point-based representations~\cite{hu2025videopredictionpolicy,liu2025geometryaware4d}, yet such view-dependent predictions struggle with multi-view consistency and do not provide the structured 3D scene information required to reason about object poses or end-effector clearances. Moreover, tightly coupling a high-capacity world model with a large VLA backbone can quickly become computationally prohibitive, especially if multi-view 3D prediction is required at inference time.

% For effective visuomotor control, we argue that two predictive capabilities are particularly valuable. First, the policy should leverage \emph{predictive kinematic priors} describing how the robot is likely to move over the next several steps, rather than relying solely on instantaneous joint states. Second, it should reason about \emph{predictive 3D scene geometry} through an explicit, differentiable representation that can be supervised using depth and multi-view signals. Crucially, such predictive signals must integrate seamlessly with VLA models while preserving deployment efficiency.


Vision--Language--Action (VLA) models have recently emerged as a powerful paradigm for robotic manipulation, leveraging large pre-trained Vision--Language Models (VLMs) to map rich visual observations and language instructions directly to actions~\cite{brohan2022rt, zitkovich2023rt, li2024cogact, zheng2025universal}. By inheriting semantic and visual priors from internet-scale data, these models generalize well across tasks and settings. However, they operate primarily in 2D image space, mapping the current observation reactively and lacking explicit 3D spatial modeling~\cite{zhong2025survey}. This 2D-centric, myopic formulation limits performance in tasks requiring precise 3D reasoning and physically consistent long-horizon control.

To move beyond purely reactive policies, recent work has begun to incorporate predictive structure into visuomotor models~\cite{hu2025video, li2025unified, lu2025gwm}. Some approaches learn latent dynamics or temporal abstractions to better capture how states evolve over time~\cite{chen2025villa, chen2025moto, li2024vision, zhao2025cot}, while others predict future observations such as RGB frames, depths, or point-based representations~\cite{cen2025worldvla,zhangdreamvla,hu2025video,qian2025wristworld, grauman2022ego4d,gupta2024pre}. Although these methods provide useful temporal signals, they often remain view-independent and do not strictly enforce multi-view or 3D geometric consistency, making it hard to reason about object poses, clearances, and end-effector motion in workspace coordinates. Moreover, tightly coupling high-capacity predictive modules with large VLA backbones can introduce non-negligible computational overhead if complex 3D prediction is required during inference time.

For robotic manipulation, we argue that two predictive capabilities are particularly valuable. First, the policy should exploit \emph{predictive kinematic priors} that summarize how the robot is likely to move over the next few steps, rather than relying solely on instantaneous joint states. Second, it should reason about \emph{predictive 3D Gaussian geometry} through an explicit, differentiable representation aligned with the robot’s workspace and amenable to supervision from depth and multi-view cues. Crucially, these predictive signals must seamlessly integrate with VLA architectures and keep inference-time overhead modest enough to enable practical deployment for real-time robotic control.

In this paper, we introduce \textbf{GeoPredict}, a geometry-aware VLA framework augmented with predictive priors for robotic manipulation. GeoPredict is built on top of a strong continuous-action VLA policy~\cite{black2410pi0}, and augments the policy with two complementary predictive modules.

First, at the kinematic level, we introduce a \emph{trajectory-level prediction module} that tracks key robot joints and end-effector points and encodes their motion history into compact tokens via a Track Encoder. A set of learnable future track queries then predicts multi-step 3D keypoint trajectories, providing explicit, differentiable kinematic priors. These predicted trajectories not only regularize the transformer's internal dynamics, but also serve as guidance for where high-fidelity 3D scene modeling is most needed.

Second, at the geometric level, we introduce a \emph{predictive 3D Gaussian geometry module} that forecasts how workspace geometry evolves over time. A coarse 3D spatial query grid is decoded into Gaussian primitives representing the scene at future timesteps. To focus geometric capacity where precision matters most, a track-guided refinement mechanism increases Gaussian density along the predicted future keypoint trajectories of the joints, yielding higher resolution around anticipated motion and interaction regions. The resulting 3DGS representation is supervised through future depth-map rendering, providing a geometry-focused training signal without requiring color modeling.

Both predictive modules are integrated via a block-wise causal attention hierarchy and are used only during training; inference behaves exactly like the base VLA policy.

In summary, our main contributions are threefold:
(1) We introduce \textbf{GeoPredict}, a geometry-aware VLA framework that injects future-aware kinematic and geometric priors into a continuous-action VLA policy, enhancing reasoning about long-horizon 3D dynamics.
(2) We propose two complementary predictive modules: a trajectory-level kinematic predictor forecasting multi-step robot keypoint motion, and a predictive 3D Gaussian geometry module with track-guided refinement that allocates geometric capacity to task-relevant interaction regions.
(3) We show that GeoPredict delivers consistent and substantial improvements over strong VLA baselines on RoboCasa Human-50, LIBERO, and real-world manipulation tasks, particularly those requiring precise spatial reasoning and geometric robustness.
